\documentclass[11pt, a4paper]{article}

\usepackage[T1]{fontenc}
\usepackage[utf8]{inputenc}
\usepackage{tgpagella}
\usepackage[spanish, es-noshorthands]{babel}
\usepackage{amsmath, amssymb, bm}
\usepackage{tabularx}
\usepackage{booktabs}
\usepackage[most]{tcolorbox}
\usepackage[margin=2.5cm]{geometry}
\usepackage{ragged2e}
\usepackage{fancyhdr}

\pagestyle{fancy}
\fancyhf{}
\setlength{\headheight}{16pt}
\lhead{\textbf{UCB Sede Tarija} | Ing. Mecatrónica}
\chead{}
\rhead{IMT-121 Circuitos Electrónicos I | Remediación Dirigida}
\lfoot{Prof: M.Sc. Ing. Bernardo Quiroga Turdera}
\rfoot{Página \thepage}
\renewcommand{\headrulewidth}{0.4pt}

\definecolor{ucbblue}{RGB}{0, 51, 102}

\begin{document}

\begin{tcolorbox}[colback=red!5!white, colframe=red!70!black, title=\textbf{Hoja de Remediación Dirigida (Path C)}]
    \centering \textbf{Corrección de Patrón de Error Dominante}
\end{tcolorbox}

\noindent\textit{\textbf{NOTA PARA EL INSTRUCTOR:} Este artefacto se encuentra deliberadamente PENDIENTE DE FINALIZAR. Según el marco metodológico del curso, la remediación (Path C) sólo se activa y se diseña UNA VEZ identificada la dificultad generalizada real de la cohorte durante la corrección del Examen Hito 2. Llenar esto con antelación sería fabricar un problema irreal.}

\vspace{0.5cm}
\section*{Estructura de la Intervención (A rellenar post-evaluación)}

\begin{enumerate}
    \item \textbf{El Error Observado (Anonimizado):} \\
    \textit{[Ej. El 60\% de la clase halló $R_{th}$ sumando todas las resistencias sin desconectar ni cortocircuitar las fuentes de tensión].}
    
    \item \textbf{Por qué falla este método:} \\
    \textit{[Ej. Si no se apaga la fuente, la red sigue inyectando energía; $R_{th}$ es por definición la impedancia de la red pasiva (muerta) vista desde los terminales externos].}
    
    \item \textbf{El Método Correcto (Paso a Paso):}
    \begin{itemize}
        \item Paso 1: \rule{10cm}{0.4pt}
        \item Paso 2: \rule{10cm}{0.4pt}
        \item Paso 3: \rule{10cm}{0.4pt}
    \end{itemize}
    
    \item \textbf{Problema de Transferencia (Reintento Independiente):} \\
    \textit{[Se debe insertar aquí un circuito nuevo, con dificultad aritmética simple y equivalente al del examen, enfocado EXCLUSIVAMENTE en la competencia fallida].}
    
    \item \textbf{Reflexión del Estudiante:} \\
    ¿Qué paso diferente aplicó esta vez que omitió durante el examen? \\
    \rule{\textwidth}{0.4pt} \\
    \rule{\textwidth}{0.4pt}
\end{enumerate}

\end{document}
